\documentclass{rapport}
\usepackage{enumitem}

\reporttitle{Vedtekter}
\projectname{Eclipse NMBU}
% \projectlogo{path/to/project-logo}   % use instead of \projectname if you have a project logo
\orglogo{../../Eclipse Images/EclipseFullText}  % path is relative to this file's folder

% ---------------------------------------------------------------------------
% Cover / title page for rapport-class documents.
%
% Usage (after \documentclass{rapport} and before \begin{document}):
%   % ---------------------------------------------------------------------------
% Cover / title page for rapport-class documents.
%
% Usage (after \documentclass{rapport} and before \begin{document}):
%   % ---------------------------------------------------------------------------
% Cover / title page for rapport-class documents.
%
% Usage (after \documentclass{rapport} and before \begin{document}):
%   \input{../cover/rapport-cover}
%   \missionpatch{path/to/patch}      % optional circular mission/team patch;
%                                     % without it, the org logo is shown large
%                                     % at the top instead
%   \coverkind{Vedtekter}             % optional small line above the name
%   \reportsubtitle{Optional subtitle line}
%   \orgname{Eclipse NMBU}{Norwegian University of Life Sciences, Norway}
%   \reportlocation{Ås, Norge}
%   \date{3. oktober 2022}            % optional, defaults to \today
%   \contactinfo{Contact: post@eclipsenmbu.no}  % optional, shown at the bottom
%
% Then, as the first thing inside \begin{document}:
%   \makecoverpage
%
% The large title line shows \projectname (falling back to \reporttitle if
% \coverkind is not set), so a document typically sets both:
%   \projectname{Eclipse NMBU}
%   \coverkind{Vedtekter}
% ---------------------------------------------------------------------------

\makeatletter
\def\@missionpatch{}
\newcommand{\missionpatch}[1]{\def\@missionpatch{#1}}

\def\@coverkind{}
\newcommand{\coverkind}[1]{\def\@coverkind{#1}}

\def\@reportsubtitle{}
\newcommand{\reportsubtitle}[1]{\def\@reportsubtitle{#1}}

\def\@orgnameone{}
\def\@orgnametwo{}
\newcommand{\orgname}[2]{\def\@orgnameone{#1}\def\@orgnametwo{#2}}

\def\@reportlocation{}
\newcommand{\reportlocation}[1]{\def\@reportlocation{#1}}

\def\@contactinfo{}
\newcommand{\contactinfo}[1]{\def\@contactinfo{#1}}

\newcommand{\makecoverpage}{%
  \begingroup
  \thispagestyle{empty}
  \begin{center}
  \ifx\@missionpatch\@empty
    \includegraphics[width=10cm]{\@orglogo}\par
  \else
    \includegraphics[width=6.5cm]{\@missionpatch}\par
  \fi
  \vspace{3.4cm}
  \ifx\@coverkind\@empty
    {\sffamily\Huge\bfseries \@reporttitle \par}
  \else
    {\sffamily\Huge\bfseries \@coverkind \par}
    {\sffamily\Huge\bfseries \@projectname \par}
  \fi
  \ifx\@reportsubtitle\@empty\else
    \vspace{1.6cm}
    {\sffamily\bfseries \@reportsubtitle \par}
  \fi
  \vspace{2cm}
  \ifx\@missionpatch\@empty\else
    \includegraphics[height=1.3cm]{\@orglogo}\par
    \vspace{0.6cm}
  \fi
  \@orgnameone \\
  \@orgnametwo \par
  \vspace{0.6cm}
  \ifx\@reportlocation\@empty
    \@date
  \else
    \@reportlocation, \@date
  \fi
  \par
  \vfill
  \ifx\@contactinfo\@empty\else
    {\small \@contactinfo \par}
  \fi
  \end{center}
  \endgroup
  \newpage
  \pagenumbering{roman}
  \pagestyle{plain}
}
\makeatother

%   \missionpatch{path/to/patch}      % optional circular mission/team patch;
%                                     % without it, the org logo is shown large
%                                     % at the top instead
%   \coverkind{Vedtekter}             % optional small line above the name
%   \reportsubtitle{Optional subtitle line}
%   \orgname{Eclipse NMBU}{Norwegian University of Life Sciences, Norway}
%   \reportlocation{Ås, Norge}
%   \date{3. oktober 2022}            % optional, defaults to \today
%   \contactinfo{Contact: post@eclipsenmbu.no}  % optional, shown at the bottom
%
% Then, as the first thing inside \begin{document}:
%   \makecoverpage
%
% The large title line shows \projectname (falling back to \reporttitle if
% \coverkind is not set), so a document typically sets both:
%   \projectname{Eclipse NMBU}
%   \coverkind{Vedtekter}
% ---------------------------------------------------------------------------

\makeatletter
\def\@missionpatch{}
\newcommand{\missionpatch}[1]{\def\@missionpatch{#1}}

\def\@coverkind{}
\newcommand{\coverkind}[1]{\def\@coverkind{#1}}

\def\@reportsubtitle{}
\newcommand{\reportsubtitle}[1]{\def\@reportsubtitle{#1}}

\def\@orgnameone{}
\def\@orgnametwo{}
\newcommand{\orgname}[2]{\def\@orgnameone{#1}\def\@orgnametwo{#2}}

\def\@reportlocation{}
\newcommand{\reportlocation}[1]{\def\@reportlocation{#1}}

\def\@contactinfo{}
\newcommand{\contactinfo}[1]{\def\@contactinfo{#1}}

\newcommand{\makecoverpage}{%
  \begingroup
  \thispagestyle{empty}
  \begin{center}
  \ifx\@missionpatch\@empty
    \includegraphics[width=10cm]{\@orglogo}\par
  \else
    \includegraphics[width=6.5cm]{\@missionpatch}\par
  \fi
  \vspace{3.4cm}
  \ifx\@coverkind\@empty
    {\sffamily\Huge\bfseries \@reporttitle \par}
  \else
    {\sffamily\Huge\bfseries \@coverkind \par}
    {\sffamily\Huge\bfseries \@projectname \par}
  \fi
  \ifx\@reportsubtitle\@empty\else
    \vspace{1.6cm}
    {\sffamily\bfseries \@reportsubtitle \par}
  \fi
  \vspace{2cm}
  \ifx\@missionpatch\@empty\else
    \includegraphics[height=1.3cm]{\@orglogo}\par
    \vspace{0.6cm}
  \fi
  \@orgnameone \\
  \@orgnametwo \par
  \vspace{0.6cm}
  \ifx\@reportlocation\@empty
    \@date
  \else
    \@reportlocation, \@date
  \fi
  \par
  \vfill
  \ifx\@contactinfo\@empty\else
    {\small \@contactinfo \par}
  \fi
  \end{center}
  \endgroup
  \newpage
  \pagenumbering{roman}
  \pagestyle{plain}
}
\makeatother

%   \missionpatch{path/to/patch}      % optional circular mission/team patch;
%                                     % without it, the org logo is shown large
%                                     % at the top instead
%   \coverkind{Vedtekter}             % optional small line above the name
%   \reportsubtitle{Optional subtitle line}
%   \orgname{Eclipse NMBU}{Norwegian University of Life Sciences, Norway}
%   \reportlocation{Ås, Norge}
%   \date{3. oktober 2022}            % optional, defaults to \today
%   \contactinfo{Contact: post@eclipsenmbu.no}  % optional, shown at the bottom
%
% Then, as the first thing inside \begin{document}:
%   \makecoverpage
%
% The large title line shows \projectname (falling back to \reporttitle if
% \coverkind is not set), so a document typically sets both:
%   \projectname{Eclipse NMBU}
%   \coverkind{Vedtekter}
% ---------------------------------------------------------------------------

\makeatletter
\def\@missionpatch{}
\newcommand{\missionpatch}[1]{\def\@missionpatch{#1}}

\def\@coverkind{}
\newcommand{\coverkind}[1]{\def\@coverkind{#1}}

\def\@reportsubtitle{}
\newcommand{\reportsubtitle}[1]{\def\@reportsubtitle{#1}}

\def\@orgnameone{}
\def\@orgnametwo{}
\newcommand{\orgname}[2]{\def\@orgnameone{#1}\def\@orgnametwo{#2}}

\def\@reportlocation{}
\newcommand{\reportlocation}[1]{\def\@reportlocation{#1}}

\def\@contactinfo{}
\newcommand{\contactinfo}[1]{\def\@contactinfo{#1}}

\newcommand{\makecoverpage}{%
  \begingroup
  \thispagestyle{empty}
  \begin{center}
  \ifx\@missionpatch\@empty
    \includegraphics[width=10cm]{\@orglogo}\par
  \else
    \includegraphics[width=6.5cm]{\@missionpatch}\par
  \fi
  \vspace{3.4cm}
  \ifx\@coverkind\@empty
    {\sffamily\Huge\bfseries \@reporttitle \par}
  \else
    {\sffamily\Huge\bfseries \@coverkind \par}
    {\sffamily\Huge\bfseries \@projectname \par}
  \fi
  \ifx\@reportsubtitle\@empty\else
    \vspace{1.6cm}
    {\sffamily\bfseries \@reportsubtitle \par}
  \fi
  \vspace{2cm}
  \ifx\@missionpatch\@empty\else
    \includegraphics[height=1.3cm]{\@orglogo}\par
    \vspace{0.6cm}
  \fi
  \@orgnameone \\
  \@orgnametwo \par
  \vspace{0.6cm}
  \ifx\@reportlocation\@empty
    \@date
  \else
    \@reportlocation, \@date
  \fi
  \par
  \vfill
  \ifx\@contactinfo\@empty\else
    {\small \@contactinfo \par}
  \fi
  \end{center}
  \endgroup
  \newpage
  \pagenumbering{roman}
  \pagestyle{plain}
}
\makeatother

\missionpatch{../../Eclipse Images/EclipsePlack}
\coverkind{Vedtekter}
\reportsubtitle{Vedtekter for Eclipse NMBU Vår 2026}
\orgname{Eclipse NMBU}{Norwegian University of Life Sciences, Norway}
\reportlocation{Ås}
\date{April 15, 2026}
\contactinfo{Contact: post@eclipsenmbu.no}

\begin{document}

\makecoverpage
\tableofcontents
\reportmainmatter

\vsection{§ 1 NAVN}

\vsubsection{§ 1.1 Foreningens navn}
Foreningens navn er Eclipse NMBU.

\vsection{§ 2 FORMÅL}

\vsubsection{§ 2.1 Formålet med Eclipse NMBU}
Eclipse NMBU er en frivillig studentforening primært for studenter som er
semesterregistrert ved Norges miljøog biovitenskapelige universitet, og skal
fungere som en sosial og faglig arena for studenter med interesse for
verdensrommet. Foreningens aktivitet skal være av frivillig karakter og rettet mot
forskning og utvikling (R\&D), med en overordnet visjon om å utvikle bærekraftig
romteknologi med utgangspunkt i FNs bærekraftsmål.

\vsection{§ 3 ORGANISASJONSFORM}

\vsubsection{§ 3.1 Status og politisk/religiøs uavhengighet}
Eclipse NMBU er en frivillig studentforening ved Norges Miljøog Biovitenskapelige
Universitet og ble stiftet 09.10.2024. Eclipse NMBU er politisk og religiøst uavhengig.

\vsubsection{§ 3.2 Juridisk person og selveid status}
Foreningen er en frittstående juridisk person med medlemmer og er selveiende. At den er
selveiende innebærer at ingen, verken medlemmer eller andre, har krav på foreningens
formue eller eiendeler, eller er ansvarlig for gjeld eller andre.

\vsection{§ 4 GENERALFORSAMLING}

\vsubsection{§ 4.1 Øverste organ}
Generalforsamlingen er foreningens øverste organ.

\vsubsection{§ 4.2 Hyppighet og tidsrammer}
Generalforsamlingen skal holdes minimum en gang i semesteret og innen utgangen av
november i høstsemesteret og utgangen av april på vårsemesteret.

\vsubsection{§ 4.3 Vedtaksførhet og stemmegivening}
Generalforsamlingen fatter vedtak ved simpelt flertall. Ingen medlemmer har mer enn én
stemme og stemmegivning kan ikke skje ved fullmakt.

\vsubsection{§ 4.4 Ekstraordinær generalforsamling}
Ekstraordinær generalforsamling kan innkalles hvis minimum 1/3 av medlemmene krever
det eller styret finner det nødvendig.

\vsubsection{§ 4.5 Innkallingsfrister og sakspapirer}
Innkalling til generalforsamling skal være medlemmene i hende minimum to uker før
generalforsamlingen. Sakspapirer skal være tilgjengelige for medlemmene minimum en uke
før generalforsamlingen.

\vsubsection{§ 4.6 Vedtektsendringer}
Vedtektsendringer som skal stemmes over på generalforsamlingen må være styret i hende
minimum en uke før generalforsamlingen og vedtektsendringer kan bare skje når 2/3 av
medlemmene stemmer for endringen.

\vsubsection{§ 4.7 Valg av styremedlemmer}
Alle styrets medlemmer skal bli valgt under generalforsamlingen og velges med flertall.

\noindent\textbf{§ 4.7.1:} Om dette ikke oppnås ved første avstemning, skal en ny avstemming
finne sted.

\noindent\textbf{§ 4.7.2:} Blanke stemmer teller som en del av stemmemassen ved votering.

Dersom ingen på valg oppnår et flertall kan blanke stemmer fjernes ut av massen ved
en ny votering. Det blir fortsatt mulig å stemme blankt, men stemmen vil ikke bli ansett
som godkjent.

\vsubsection{§ 4.8 Saksliste}
Sakliste for generalforsamlingen skal følge faste punkter:
\begin{enumerate}[label=\alph*.]
  \item Valg av ordstyrer.
  \item Valg av referent.
  \item Valg av nøytralt tellekorps: to medlemmer.
  \item Valg av protokollunderskrivere.
  \item Godkjenning av regnskap fra forrige semester.
  \item Godkjenning av regnskap for kommende semester.
  \item Semesterberetning / styreleders beretning.
  \item Valg av styremedlemmer.
  \item Eventuelt annet.
\end{enumerate}

\vsection{§ 5 STYRET}

\vsubsection{§ 5.1 Styrets sammensetning}
Styret sammensetning skal være følgende:
\begin{itemize}
  \item Styreleder
  \item Økonomi ansvarlig
  \item Teknisk ansvarlig
  \item Markedsførings ansvarlig
\end{itemize}

\vsubsection{§ 5.2 Styreleder (CEO: Chief Executive Officer)}
\begin{itemize}
  \item \textbf{Ansvar:} Overordna administrativ ledelse, ekstern representasjon, kontraktssignering
        og strategisk planlegging. CEO har det øverste ansvaret for at organisasjonen
        opererer innenfor norsk lov og universitetets reglement. CEO har videre ansvar for
        organisasjonskultur, HMS-protokoller og overordna koordinering mellom de ulike
        lederrollene.
  \item \textbf{Fullmakt:} CEO har avgjerande stemme ved stemme-likhet i administrative
        styresaker. CEO har fullmakt til å representere organisasjonen utad og inngå
        bindende avtaler på vegne av foreiningen innenfor rammene sett av styret.
  \item \textbf{Varighet:} CEO er et verv som varer i 2 student semester.
\end{itemize}

\vsubsection{§ 5.3 Teknisk Ansvarlig (CTO: Chief Technology Officer)}
\begin{itemize}
  \item \textbf{Ansvar:} Ansvarlig for alle tekniske prosjekt. Ansvarlig for systemintegrasjon, teknisk
        sikkerhet og all teknisk framdrift. CTO skal sørge for at den tekniske utviklingen
        følger fastsatte tidsplaner og milepæler, og har ansvar for ressursstyring og logistikk
        knyttet til teknisk utstyr og verkstad.
  \item \textbf{Fullmakt:} CTO har fullmakt i tekniske beslutninger som berører systemets integritet
        eller sikkerheit. CTO har fullmakt til å kalle inn til tekniske statusmøte.
  \item \textbf{Varighet:} CTO er et verv som varer i 2 student semester.
\end{itemize}

\vsubsection{§ 5.4 Økonomi ansvarlig (CFO: Chief Financial Officer)}
\begin{itemize}
  \item \textbf{Ansvar:} Finansiell planlegging, bokføring, budsjettkontroll og revisjon. CFO skal syte
        for at foreininga har en sunn økonomi og at midler blir brukt i tråd med vedtekne
        budsjett. CFO har ansvar for forsikring og finansiell risikohåndtering.
  \item \textbf{Fullmakt:} CFO har fullmakt til å fryse utbetalinger til et prosjekt dersom
        budsjettrammene blir overstridet eller nødvendig dokumentasjon mangler. CFO skal
        godkjenne alle større innkjøp før bestilling.
  \item \textbf{Varighet:} CFO er et verv som varer i 2 student semester.
\end{itemize}

\vsubsection{§ 5.5 Marknadsførings ansvarlig (CMO: Chief Marketing Officer)}
\begin{itemize}
  \item \textbf{Ansvar:} Rekruttering, merkevarebygging, sponsorrelasjon og ekstern
        kommunikasjon. CMO skal sørge for at foreininga er synlig og attraktiv for både nye
        medlemmer og samarbeidspartnere, og har ansvar for innhald og vedlikehald av
        organisasjonen sine digitale flater.
  \item \textbf{Fullmakt:} CMO godkjenner alt materiale som blir publisert eksternt (sosiale medier,
        nettside, pressemeldinger) og leder organisasjonen sine rekrutteringskampanjer.
  \item \textbf{Varighet:} CMO er et verv som varer i 2 student semester.
\end{itemize}

\vsubsection{§ 5.6 Konstituering av fratreden}
Dersom et styremedlem har behov for å fratre sin stilling kan styret konstituere et nytt
midlertidig medlem ved alminnelig flertall inntil stillingen fylles ved generalforsamling. Det
konstituerte medlemmet har stemmerett i styremøter.

\vsubsection{§ 5.7 Erfaringsoverføring}
Når et nytt styremedlem inntrer, skal det forrige styremedlemmet sørge for en god
erfaringsoverføring. Dette innebærer at det gamle styremedlemmet skal være tilgjengelig for
styret og nytt styremedlem i minst 90 dager etter det gamle styremedlemmet avtrer.

\vsubsection{§ 5.8 Stemmegivning i styret}
Ved likt antall stemmer ved saker innad i styret teller styreleder sin stemme dobbelt.

\vsection{§ 6 START AV PROSJEKTER}

\vsubsection{§ 6.1 Beslutningsprosess for prosjekter}
Alle større prosjekter som igangsettes av Eclipse NMBU, skal igangstartes og
stemmes over av styret. Dersom medlemmer har forslag til prosjekter som de vil
gjennomføre kan dette sendes inn til styret, så vil styret ta en vurdering på om
det er gjennomførbart eller ikke.

\vsection{§ 7 ETIKK OG MILJØ}

\vsubsection{§ 7.1 Bærekraft og etiske hensyn}
Foreningen skal så langt det er mulig vurdere miljømessige og etiske hensyn i planleggingen
av aktiviteter. Foreningen skal i møte med andre organisasjoner, foreninger, arrangementer
og oppdragsgivere etterspørre alternativer slik at etikk og miljø blir tatt hensyn til.
Foreningen skal ta hensyn til FNs bærekraftsmål og bidra så langt det lar seg gjøre å
imøtekomme målene. Oppbyggingen, struktur og organisering skal foregå med den ånd at
foreningen er bærekraftig på en sånn måte at kommende generasjoner har anledning til å
bygge videre på tidligere aktiviteter.

\vsection{§ 8 VARSLING OG HÅNDTERING AV KRITIKKVERDIG FORHOLD}

\vsubsection{§ 8.1 Definisjon av varsling}
Varsling er at et medlem, et styremedlem eller noen som er i kontakt med foreningen sier
ifra om et kritikkverdig forhold i forbindelse med foreningens aktiviteter.
Eksempler på kritikkverdige forhold kan være
\begin{itemize}
  \item Helse, miljø eller sikkerhetsmessige forhold
  \item Personkonflikter
  \item Forstyrrende adferd
  \item Lovbrudd som tyveri, økonomisk utroskap eller underslag
  \item Seksuell trakassering
  \item Overgrep
\end{itemize}

\vsubsection{§ 8.2 Lovbrudd og saksbehandling}
Ved alvorlige beskyldninger om lovbrudd må forhold anmeldes til politiet for at styret skal
kunne behandle varselet. I saker der varselet er rettet mot en eller flere personer skal alle
involverte få anledning til å kunne gi sin versjon av situasjonen. Lovbrudd og seksuell
trakassering er ikke akseptert og styret kan anmelde forhold basert på egne observasjoner.

\vsubsection{§ 8.3 Vedtak om reaksjoner}
Vedtak om reaksjoner gjøres med 2/3 stemming i foreningens styre.

\vsubsection{§ 8.4 Habilitet og sanksjoner}
Ved varsling av kritikkverdige forhold skal det vurderes hvorvidt styret for en slik håndtering
er habile. Dersom utnevnte viser seg å være inhabil i håndtering av varslingssak skal et
uavhengig organ oppnevnes.

Reaksjoner etter varsel:
\begin{itemize}
  \item Endringer i foreningens rutiner
  \item Advarsel
  \item Suspensjon/utestengelse fra styre/verv og foreningens aktiviteter for en periode
  \item Permanent utestengelse fra styre/verv og foreningens aktiviteter.
\end{itemize}

\vsection{§ 9 PERSONVERN}

\vsubsection{§ 9.1 Behandling av personopplysninger}
Personvern for medlemmer skal bli ivaretatt ved at det bare behandles opplysninger som er
nødvendige og relevante for å administrere medlemskapet. Ved et opphør av medlemskap
skal medlemmets opplysninger slettes uten oppfordring fra det opphørte medlemmet, og
skal være i linje med datatilsynets retningslinjer. Medlemmet skal motta bekreftelse på
sletting uten ugrunnet opphold.

\vsection{§ 10 RETTIGHETER OG PLIKTER KNYTTET TIL MEDLEMSKAPET}

\vsubsection{§ 10.1 Definisjon av medlemsnivåer}
Eclipse NMBU opererer med fire medlemsnivåer basert på dokumentert aktivitet og ansvar i
inneværende semester. Medlemskap varer i ett år (to semestre) om ikke annet er avtalt, men
aktivitetsnivået vurderes per semester. Nivåene er hierarkiske; høyere nivåer innehar alle
rettigheter og plikter tilhørende nivåene under.
\begin{enumerate}
  \item \textbf{Styremedlem:} Innehar det juridiske og administrative ansvaret for foreningen.
  \item \textbf{Kjernemedlem:} Medlemmer med utvidet ansvar og mulighet til lederansvar for
        spesifikke oppgaver/undergrupper.
  \item \textbf{Aktivt medlem:} Medlemmer som oppfyller fastsatte krav til poengopptjening.
  \item \textbf{Inaktivt medlem:} Medlemmer som er registrert i foreningen, men som ikke oppfyller
        kravene til aktiv deltakelse.
\end{enumerate}

\vsubsection{§ 10.2 Rettigheter og stemmerett}
\begin{enumerate}
  \item \textbf{Stemmerett og valgbarhet:} Alle Styremedlemmer, Kjernemedlemmer og Aktive
        medlemmer har stemmerett ved generalforsamlingen og er valgbare til tillitsverv i
        foreningen.
  \item \textbf{Møterett:} Alle medlemmer, inkludert inaktive, har rett til å delta på
        generalforsamlingen.
  \item \textbf{Inaktive medlemmer:} Inaktive medlemmer mangler stemmerett og valgbarhet. Ved
        overgang til inaktivt medlemskap fratas stemmeretten sammen med andre tilhørende
        rettigheter.
\end{enumerate}

\vsubsection{§ 10.3 Medlemmenes plikter}
Medlemmene plikter å forholde seg til vedtak fattet av generalforsamlingen og styret. For å
sikre progresjon gjelder følgende:
\begin{enumerate}
  \item \textbf{Leveringsplikt:} Medlemmer som har påtatt seg en oppgave, plikter å levere denne
        innen gitt tidsramme.
  \item \textbf{Delegasjonsplikt:} Hvis et medlem ikke kan gjennomføre en påtatt oppgave, plikter
        medlemmet å delegere denne videre eller varsle ansvarlig
        styremedlem/prosjektleder.
  \item \textbf{Fravær:} Dersom et medlem er inaktivt uten gyldig grunn over en sammenhengende
        periode på 4 uker i en aktiv arbeidsperiode, kan styret vurdere opphør av
        medlemskapet.
        \begin{enumerate}[label=\alph*.]
          \item Grunn til fravær skal avklares med medlemmets nærmeste leder i forhold til
                Eclipse sin styringsstruktur. Aktiv periode følger NMBUs definisjon av vårog
                høstparallell (oppstart til eksamensperiodens begynnelse).
        \end{enumerate}
\end{enumerate}

\vsubsection{§ 10.4 Poengsystem}
Eclipse NMBU benytter et poengbasert system for å måle aktivitet og verdiskaping.
\begin{enumerate}
  \item \textbf{Tildeling:} Poeng tildeles basert på utførte og godkjente leveranser, teknisk
        progresjon, oppmøte og administrativ innsats.
  \item \textbf{Forvaltning:} Styret fastsetter og annonserer poengsatser og grenser for hele
        foreningen før hver periode.
  \item \textbf{Kvalitetssikring:} Godkjennelse av poenggivende arbeid gjøres av et styremedlem
        eller oppgaveansvarlig.
\end{enumerate}

\vsubsection{§ 10.5 Styrets fullmakter vedrørende medlemskap}
Utenom generalforsamlingen kan styret ved flertallsvedtak foreta følgende endringer:
\begin{enumerate}
  \item \textbf{Nivåendring:} Endre medlemsnivået til et medlem basert på poengopptjening eller
        manglende overholdelse av plikter.
  \item \textbf{Verv:} Endre hvem som besitter et verv, samt tildele eller avslå søknader om verv og
        medlemskap.
  \item \textbf{Eksklusjon:} Fjerne medlemmer fra foreningens medlemsliste dersom medlemskapet
        ikke er fornyet, eller dersom vilkårene i § 10.3 punkt 3 er brutt.
\end{enumerate}

\vsection{§ 11 SIGNATURRETT}

\vsubsection{§ 11.1 Utøvelse av signaturrett}
Det er styrets medlemmer som i fellesskap innehar signaturrett. Dersom hele styret har i
fellesskap kommet til enighet om å signere er det tilstrekkelig at enten styrets leder eller
økonomi ansvarlig utfører den faktiske signeringen. Det sittende styret står ansvarlig for alle
gjeldende kontrakter.

\vsection{§ 12 ØKONOMI}

\vsubsection{§ 12.1 Bruk av midler}
Foreningens midler skal kun brukes i henhold til foreningens formål.

\vsubsection{§ 12.2 Budsjett og regnskap}
Det skal alltid foreligge budsjett og regnskap for foreningens midler, med tilhørende bilag.
Budsjett og regnskap gjennomgås og godkjennes på generalforsamling.

\vsubsection{§ 12.3 Godkjenning av utbetalinger}
Utbetalinger skal alltid godkjennes av leder: eller leders stedfortreder: i tillegg til
økonomiansvarlig.

\vsubsection{§ 12.4 Økonomiske investeringer}
Ved stemming over økonomiske investeringer trengs det et 2/3 flertall i styret.

\vsubsection{§ 12.5 Økonomisk ansvar}
Økonomisjef har ansvar for oppfølging av alle økonomiske saker og plikter å
informere styret fortløpende om den økonomiske bruken til organisasjonen.

\vsection{§ 13 OPPHØR}

\vsubsection{§ 13.1 Prosedyrer for oppløsning}
Foreningen kan oppløses ved at 2/3 av de stemmeberettigede medlemmene på
generalforsamlingen stemmer for oppløsning. Sak om oppløsning må være meldt inn til
ordinære saksfrister i forkant av generalforsamlingen. Ved opphør vil foreningens midler og
eiendeler bli gitt til en annen egnet studentforening med samme formål, som
generalforsamlingen bestemmer.

\vsubsection{§ 13.2 Disponering av eiendeler}
Ingen medlemmer har krav på foreningens midler eller andel av disse.

\end{document}
